\chapter{The Machine}\label{cha:machine}

The K4510 is a fantasy computer: a machine that never existed, built the
way 1985 might have built it with no budget committee in the room. It is not
an emulation of any real computer. Its CPU, video chip, sound, and operating
software are its own; the parts it borrows --- a 6502-family instruction set,
the AdLib's FM chip --- it borrows openly and then outgrows.

\section{What is the K4510?}
\begin{itemize}[leftmargin=1.4em]
\item \textbf{CPU: the 45GS10} --- the MEGA65's 45GS02 instruction set (a
  4510 with the Q pseudo-register and 32-bit flat addressing) plus this
  machine's own memory management: bank registers, a far-call gate, and RAM
  under the ROM. Its clock is whatever the host it runs on can hold at
  sixty frames a second with no gaps in the sound: the machine is a
  fantasy and its timings are suggestions. The clock is a setting (F7,
  under Machine --- six steps from 60\,MHz down to 10; the ladder goes on
  to 202.5, and those faster steps are held back for now). A host nobody
  has measured runs at 40.5, the MEGA65's number; \verb|SETUP| measures
  it properly, with the sound and the picture really running, and the
  machine keeps the answer for that host so later boots pay nothing. It
  never second-guesses a clock you chose yourself. A desktop of the last
  ten years holds a hundred megahertz or more. \verb|INFO| reports the
  clock in force, read live, and \verb|BENCH| is there when you want the
  whole picture.
\item \textbf{Memory: 256\,MB}, flat, 28-bit. The CPU sees 64\,KB at a time;
  everything else is one instruction away.
\item \textbf{Video: VICKY} --- 640$\times$480, 256 colours from a 24-bit
  palette, four layers, 128 sprites, a blitter that draws lines and filled
  triangles, and a display-list coprocessor named SHEILA. Smaller modes
  (640$\times$240, 320$\times$240, 320$\times$200, 160$\times$200) are
  drawn into the same glass, doubled and centred, so the raster is always
  480 lines however few of them the picture uses.
\item \textbf{Sound: MELODY} --- an OPL2, the Yamaha YM3812, at
  \texttt{\$D480}: nine FM voices wired the AdLib's way, an address port,
  a data port and a status register you poll. Any AdLib register list or
  instrument patch therefore means what it says on this machine. A
  four-channel sound sequencer at \texttt{\$D5E0} plays through it, in the
  BBC Micro's idiom, which is what BBC BASIC's \verb|SOUND| and Mad
  Pascal's \verb|Sound| use. A floating-point MATH unit sits beside it,
  which the BASICs and LOGO lean on. (The machine had four SID chips
  until September 2026. Appendix~\ref{cha:sound} is the whole story.)
\item \textbf{K/OS}: a shell with directories, in ROM --- and on the disk,
  one word away each, seven languages: EhBASIC with graphics and sprites,
  Microsoft's own BASIC of 1977, BBC BASIC on the Tube, Forth, LOGO, CP/M
  on a Z80, and RX, the machine's REXX. Two compilers are one word away
  as well: \verb|PAS| and \verb|CC| turn a Pascal or C source in the
  directory you are standing in into a program beside it
  (Chapter~\ref{cha:linux}).
\end{itemize}

\section{One machine, three ways to run it}\label{sec:names}
The computer is the \textbf{K4510}: the 45GS10, VICKY, SHEILA, MELODY
and K/OS. There is one of it. What changes is what it is standing on.

\textbf{As a whole computer, from a stick.} Write the image
(\pth{linux/build-live.sh}) to a USB stick, boot a spare laptop from it,
and the machine is the machine: a minimal Debian that exists only to
hold it up, loaded into RAM, keeping its settings and your files on the
stick and never touching the laptop's own drive.

\textbf{As a whole computer, installed beside another system.}
\pth{linux/install-k4510.sh} copies the same system from the stick onto a
partition of a computer's internal disk and adds it to that computer's
boot menu, so the machine is one choice of two when it starts. The other
system is not changed, and the computer boots whichever of the two was
used last (Chapter~\ref{cha:linux}).

Either way, the Linux underneath is not hidden. The cross-compilers,
git, an editor and a second terminal are on it, and the machine can
reach them (Chapter~\ref{cha:linux}); but it is furniture. Switch on and
you are at the \verb|/HOME]| prompt.

\textbf{As a window.} Run the same program on a Linux desktop and you
get the same machine in a window, which is how it is developed and
tested, and how most of the figures in this book were taken. There is
also a sandboxed container flavour (\pth{linux/podman.sh}) that is shown
nothing of your computer but the display, the sound, the game
controllers and one shared folder.

Same ROM bytes, same software, every way. Where the desktop window and
the whole computer differ, the book says ``on a desktop'' or ``on the
K4510's own Linux''. (Until 2026-09-07 there was also a bare-metal
edition for the Raspberry~Pi~3B+; it was retired so that every feature
is built once. A Pi still runs the machine under Linux like any other
computer; the tag \texttt{alpha-0.5} is the last tree with the bare-metal
edition in it.)

\section{Building the machine}
On a \textbf{Linux desktop}, three lines build the whole machine:
\begin{type}
git clone https://github.com/mlongval/k4510
cd k4510 && ./setup.sh
./k4510
\end{type}
\verb|setup.sh| installs the compilers and SDL2, builds everything, and
runs the machine's test suites.

\subsection*{Or: the whole computer}
\begin{type}
sudo ./linux/build-live.sh
\end{type}
\noindent makes the bootable image: a small Debian, the machine and its
compilers, about half an hour the first time. Write it to a USB stick
with \texttt{dd} and boot from it; to put it on a computer's own disk
instead, boot that computer's usual Linux with the stick in and run
\verb|sudo ./install-k4510.sh| from it. Chapter~\ref{cha:linux} has both.

\subsection*{What it needs}
The machine asks little of a computer, and most of what it asks is for
the Linux under it:
\begin{center}
\small
\begin{tabular}{@{}l >{\raggedright\arraybackslash}p{0.66\linewidth}@{}}
\toprule
Processor & 64-bit x86, two cores; a Core~2 is enough to run it, and what the processor decides is only the clock \verb|SETUP| settles on \\
Memory & 2\,GB; 4 is comfortable. The whole system is loaded into RAM at boot (about 0.9\,GB), and a running machine uses 1.7\,GB in all, the emulator itself 90\,MB of it \\
Graphics & any Intel or AMD chip Linux drives natively (KMS); no desktop is used \\
Screen & 640$\times$480 or more \\
Firmware & UEFI or BIOS, either \\
Storage & a 2\,GB USB stick, or a 2\,GB partition \\
\bottomrule
\end{tabular}
\end{center}
\noindent In a window on a desktop it needs only what the desktop has
already: SDL2 and a few hundred megabytes of memory.

\section{First boot}
Switch it on (on a desktop: \verb|./k4510|; otherwise, power up). The
machine's logo is on the screen while Linux starts, and then the machine
greets you in yellow on blue:

\screen{shots/boot.png}{The boot screen: five stepped colour bars, then
the machine, its CPU, its memory and its chips --- and no clock, because
the machine no longer has one speed. Everything else, the clock in force
included, is one \texttt{INFO} away, and \texttt{BANNER} prints this again.}

The \verb|/HOME]| at the bottom is the shell prompt --- the part before
the \verb|]| is the directory you are in, and the machine starts you in
\pth{/HOME}, which is yours. Type \verb|HELP| for the commands,
\verb|INFO| for the machine's full self-description, and \verb|DIR| to
see your files. The first time on a new computer, type \verb|SETUP| and
let the machine measure it.

\begin{aside}
\textbf{Keys:} Escape is RUN/STOP (stops a BASIC program), Ctrl-C is
STOP, Shift+Escape quits the emulator. Reset is a chord, the way it was
Commodore+Restore on the C64: Super+PageUp (the Windows or Command key
and PageUp) --- one key cannot do it by accident. F7 opens the settings
menu; F8 pauses the machine and turns the keyboard into a debugger's.
PrtSc takes a screenshot --- the screen flashes, and the picture is in
\pth{shots/} beside the emulator.
\end{aside}

\section{The keyboard}
The keyboard is the one in front of you, and F7~$\rightarrow$~Input
$\rightarrow$~\emph{Keyboard layout} says what its keys mean: US,
US-International, Canadian French, French, German, Spanish, UK or
Italian --- or \emph{Host}, the first choice, which leaves the question
to the desktop or Linux the machine is running on. A layout chosen here
takes effect at once, with its AltGr and its dead keys: on US-intl,
\verb|'| then \verb|e| is \'e, the dead key twice or before a space
gives the mark itself, and a \verb|'| or \verb|"| before a letter that
takes no accent types itself. On the K4510's own Linux the Linux consoles
follow the same choice. A letter the machine's character set (the IBM
PC's) cannot draw is not offered. For programming, plain US has no dead
keys to step round.

\emph{Caps Lock is Ctrl}, beside it, makes the Caps Lock key a second
Ctrl, where the old Unix keyboards had it; the shell's own
\verb|CAPSLOCK| command (Chapter~\ref{cha:shell}) is how you get capitals
without it.

\section{The F7 menu}
F7 (unshifted; Shift+F7 still reaches BBC BASIC) freezes the machine and
takes the whole screen, in the manner of the C64~Ultimate's: the
categories down the left, the settings of the chosen one on the right,
and a line at the foot saying what the keys do in whichever pane holds
the cursor. Sound stops, and closing the menu resumes exactly where it
froze. The menu draws with its own copy of the font and never touches the machine's
memory, so it opens even when a program has wrecked the screen.

Up and Down move; Enter or Right crosses from the categories to the
settings; Left and Right step a value; Escape goes back a pane, and
closes at the categories; F7 closes from anywhere. The mouse works too.
A setting with a list of choices opens that list over the panes, and
\emph{applies as the cursor passes each one} --- the scaling changes
under the menu so you can see it --- with Escape putting back whatever
was there when the list opened.

\screen{shots/menu.png}{The F7 menu: categories on the left, the Video
page on the right.}

\begin{description}
\item[Video] the border; \emph{resolution}, which is the
  \verb|MODE| command's knob;
  \emph{scaling}: \emph{integer}, to begin with (a whole-number
  multiple, so every pixel is the same size on the glass, with a border
  where the window is not an exact multiple), or \emph{fit to display}
  (as large as the window takes; the pixels may come out uneven) ---
  both hard pixels, never blurred; full screen; \emph{vertical sync}, off to begin
  with --- see the aside below; and \emph{placement} and \emph{side
  panel}, which are the section after next. The figures in this book are
  taken with the effects off.
\item[Terminal] \emph{status bands}, on or off (a row at the top, a row at the bottom)
  (Section~\ref{sec:bands}), and the clock and date format they print.
\item[Audio] volume.
\item[Input] the reset chord; which key opens the menu; whether a click
  captures the mouse pointer, and whether the host's pointer shows over
  the picture (full screen, the pointer stays on the machine's picture,
  and goes into the side panel only when there is one); the keyboard, above; and the \emph{key pipe}, typing from
  another computer (Chapter~\ref{cha:linux}): \emph{off}, \emph{on}, or
  \emph{on, shown} --- the one it starts at --- where every key typed
  that way is echoed in a bar at the foot of the window for a few
  seconds, so nobody types into the machine unseen.
\item[Machine] \emph{save state} and \emph{load state}, four slots each
  (the whole machine --- CPU, every used page of the 256~MB, VICKY, the
  devices, JIM --- to \texttt{k4510-slotN.k4s} beside the settings file;
  the Tube co-processor is not in the file and is stopped by a load);
  reset; power cycle; stop the Tube; quit; \emph{CPU clock} --- the
  steps, live, 60\,MHz at the most for now; choosing one switches \emph{Auto clock} off, because a
  clock chosen by hand is not to be second-guessed. \emph{Auto clock} on
  uses what \verb|SETUP| measured on this host. On the K4510's own Linux
  there is a last row, \emph{Shut down the computer}.
\item[Shell] whether an unknown word at the prompt may run a CP/M
  \texttt{.COM} (Chapter~\ref{cha:cpm}) --- off to begin with, on
  purpose: \verb|D| typed for \verb|DIR| should not start a Z80 program;
  and whether \texttt{/STARTUP.BAT} runs at power-on, which is the way
  out of a bad one (Section~\ref{sec:badstartup}).
\item[Info] the version and the exact build, the ROM, the files, the
  host --- the K4510's own Linux, or a window on a desktop --- and the
  battery, where there is one.
\item[Host] on the K4510's own Linux only: the computer's name, its
  network address and Tailscale's; \emph{Lid closed}, which is
  \emph{keep running} to begin with --- the machine goes on with the lid
  down --- or \emph{suspend}; a \emph{Wi-Fi / network setup} row that
  opens the network settings on a spare console, and a row that telnets
  into the Linux for you.
\end{description}

Leaving the menu writes the settings to \texttt{k4510.cfg} beside
\texttt{fs/} --- a plain \texttt{key = value} file you may edit;
comments and keys it does not know are kept, which is how a settings
file survives the machine changing under it.

\begin{aside}
\textbf{Vertical sync is off, and that is deliberate.} With it off the
machine keeps its own 60\,Hz and presents a frame when it is ready; the
cost is tearing. With it on, the pacing belongs to your display, and on
a host whose refresh is not exactly 60 that costs frames --- and frames
are sound here (Appendix~\ref{cha:sound}). Neither answer is right for
every host, which is why it is a row and not a decision.
\end{aside}

\section{The menu file: what the menu shows}\label{sec:menufile}
Beside \texttt{k4510.cfg} is a second file, \texttt{k4510-menu.cfg}, and
it decides not what the settings \emph{are} but which of them the menu
\emph{offers}. The machine writes it out in full the first time it
starts --- every category and every row, each marked \verb|show| --- so
the file is itself the list of what can be changed:

\begin{form}
[K4510]
Audio                    = show
...
[Video]
Border width             = hide
...
[Locks]
linux    = open
consoles = open
\end{form}

\begin{itemize}[leftmargin=1.4em]
\item \verb|hide| on a row takes it out of the menu, and a category with
  nothing left in it goes too; \verb|hide| under \verb|[K4510]| takes a
  whole category away. A hidden setting keeps the value it has in
  \texttt{k4510.cfg}: it is not reset, only out of reach.
\item \verb|linux = locked| shuts every door into the Linux
  underneath: \verb|!|, \verb|SSH|, and the telnet row of the Host
  menu. \verb|PAS|
  and \verb|CC| still compile. \verb|consoles = locked| makes
  Ctrl+Alt+F2 to F6 do nothing.
\item Names are the menu's own, in either case; \verb|#| starts a
  comment; a row the file does not name is shown. The file is read when
  the machine starts, so a change takes effect at the next start.
\end{itemize}

\noindent The file lives outside the machine's disk, so nothing running
on the machine can undo it. On the K4510's own Linux it is in
\pth{/home/k4510/k4510/}; once Linux is locked, edit it over ssh from
another computer.

\subsection*{Every row}
The list below is generated from the menu's own source, with the name
each setting has in \texttt{k4510.cfg}, its choices and where it starts.

\input{generated/menu}

\section{Placement, the side panel, and F8}
A 4:3 picture on a 16:9 screen leaves a third of the glass empty.
\emph{Placement} (centre, left, right) puts the picture against one edge
instead of in the middle, and \emph{side panel} fills what is left ---
at present with one thing, the machine as the emulator sees it. A panel
with a centred picture makes no sense, so turning the panel on forces
the picture left.

The panel is a strip twenty-six columns wide taking the whole height of
the window, one fact to a line: the program counter, A~X~Y~Z, the stack
pointer, the B register, the flags, the next few instructions
disassembled, VICKY's mode and the raster line it is on, the eight bank
registers with the engaged ones lit, the audio gaps, the frame counter
and the frame rate. It never reads through the I/O page --- a read of
\texttt{\$D100} would pop a key off the keyboard, and an instrument that
changes what it measures is not one.

\textbf{F8 pauses the machine, and then the keyboard is the debugger's:}

\begin{center}
\small
\begin{tabular}{@{}l >{\raggedright\arraybackslash}p{0.60\linewidth}@{}}
\toprule
\textbf{Key} & \textbf{Does} \\
\midrule
F8 & pause, and again to run on \\
Space & one instruction \\
\texttt{L} & the rest of this scanline \\
\texttt{F} & the rest of this frame \\
\texttt{D} & write a dump --- \pth{dumps/dump-NNN.txt}, the same file \texttt{DUMP} writes, and arm the instruction recorder so the next one has history \\
\texttt{T} & start or stop an instruction trace to \pth{/SYSTEM/LOG/TRACE.TXT}: one line per instruction with the registers after it, stopping itself at 200{,}000 lines \\
\bottomrule
\end{tabular}
\end{center}

\noindent Every other key is swallowed while the machine is paused. The
panel shows the legend and the state: which scanline is next, whether the
trace is running and how long it is, and the number of the last dump.
